\section{Pregunta principal y objetivos}
\label{sec:pregunta_objetivos}

% Formular la pregunta central: ¿Las variables latentes del modelo nnELM
% tienen correlatos en la señal EEG registrada durante una tarea de
% búsqueda visual híbrida?
%
% Enunciar los tres objetivos específicos:
% 1. Extraer y validar variables del modelo nnELM sobre el dataset HSEM.
% 2. Asignar esas variables a las fijaciones reales del pipeline de EEG
%    de Damián (problema de celdas vs. píxeles y sincronización).
% 3. Incorporar las variables al modelo de regresión EEG y evaluar su
%    contribución con criterios de selección de modelos.

El modelo \textit{nnELM} genera, fijación por fijación, un conjunto de mapas que representan distintos aspectos del proceso de búsqueda: el mapa de saliencia espacial, el mapa de similitud con el target, el mapa de visibilidad periférica, y el posterior bayesiano resultante de integrar evidencia de forma acumulativa a lo largo del scanpath. Por otro lado, el framework de TRFs por regresión regularizada permite estimar, a partir de la señal EEG continua registrada con los movimientos oculares, la contribución de predictores definidos a nivel de fijación a la respuesta neural. Esto abre una pregunta que articula ambas líneas de investigación:

\begin{tcolorbox}[
    colback=gray!10,
    colframe=black!60,
    boxrule=0.8pt,
    arc=4pt,
    left=10pt, right=10pt, top=8pt, bottom=8pt
]
\textbf{Pregunta Principal:}{ Dado que contamos con los mapas generados por el modelo nnELM y con el framework de análisis de TRFs, ¿es posible diseñar y extraer variables a nivel de fijación a partir de dichos mapas que tengan correlatos en la señal EEG registrada durante una tarea de búsqueda visual híbrida?}
\end{tcolorbox}

\noindent Esta pregunta es relevante desde dos puntos de vista. Por un lado, desde una perspectiva teórica, permitiría identificar qué aspectos del proceso de búsqueda, definidos computacionalmente por el modelo, se manifiestan en la actividad neural. Por otro lado, desde una perspectiva metodológica, implica compatibilizar dos pipelines desarrollados de forma independiente, con distintos criterios de preprocesamiento y distintas resoluciones espaciales: el modelo nnELM opera sobre una representación discreta de la escena en forma de grilla de celdas, mientras que el algoritmo de detección de fijaciones del pipeline de EEG opera a nivel de píxel.

\medskip
\noindent Para abordar la pregunta, se plantean tres objetivos específicos:

\begin{enumerate}

    \item \textbf{Diseño, extracción y validación de variables del modelo \textit{nnELM} sobre el dataset HSEM.}
    Diseñar variables a nivel de fijación a partir de los mapas generados por el modelo, aplicarlo sobre los scanpaths del dataset HSEM y extraer dichas variables para cada fijación. Validar su distribución, coherencia con las condiciones experimentales y su comportamiento a lo largo de la progresión del trial.

    \item \textbf{Compatibilización de los pipelines de movimientos oculares y EEG.}
    Dado que el modelo \textit{nnELM} opera sobre la representación en celdas de la escena, y que el pipeline de EEG utiliza los timestamps de fijaciones detectados por un algoritmo distinto, es necesario resolver el problema de alineación entre ambos conjuntos de fijaciones. Esto incluye definir una estrategia de asignación de las variables del modelo a las fijaciones del análisis EEG, y analizar los casos en que múltiples fijaciones consecutivas corresponden a la misma celda del modelo.

    \item \textbf{Evaluación de la contribución de las variables a la predicción de la señal EEG.}
    Incorporar las variables del \textit{nnELM} como predictores adicionales en el modelo de regresión regularizada, siguiendo el esquema de crecimiento gradual de complejidad de \citet{care2025}. Evaluar su contribución mediante $\Delta$AIC con grados de libertad efectivos de Ridge, diagnosticar la multicolinealidad con VIF, y determinar la significancia estadística de los efectos mediante TFCE.

\end{enumerate}